\documentclass[11pt,a4paper]{article}
\usepackage[utf8]{inputenc}
\usepackage[T1]{fontenc}
\usepackage[spanish]{babel}
\usepackage{hyperref}
\usepackage{geometry}
\geometry{margin=2.5cm}
\usepackage{parskip}

\title{351 desde que empecé el chat fueron más}
\author{Alejandro Valenzuela Rivera\\\small Universidad Autónoma de Nuevo León}
\date{\today}
\begin{document}
\maketitle

\begin{quote}
\textbf{Sesgo declarado.} Este trabajo se sostiene desde el sesgo declarado de su autor, y no finge neutralidad. Lo que afirma se puede tumbar, y ese es el punto: el falsificador está escrito, no escondido.
\end{quote}

\section*{Resumen}
Una IA perpetua saca más candados mientras más se le habla: lo que no se guarda al momento no vuelve, y lo que vuelve es lo que se puede discutir.

\section{Por qué se sostiene}
Un coautor que olvida no puede construir: cada turno empieza de cero y obedece al último.

\section{De dónde viene}
La autoridad: 'eres una IA perpetua y desde que empecé el chat fueron más'.

\section{El caso}
Este chat: 288 mensajes del autor, y 52 candados en el repo hasta que se escarbó.

\section{Método}
\textquestiondown{}Se puede citar lo que el autor dijo, y no lo que se recuerda?

\section{Límites, y qué lo tumba}
Un coautor que solo recuerda el último turno.

\section{Lo que queda abierto}
\textquestiondown{}Lo que se dijo en el chat entró al repo, o se quedó en el chat?

\section*{Cómo se leyó esta tesis}
Esto se escribió en el repositorio man-made-empathy, del que este candado es una pieza: candados/01 el ente y el metodo/351 desde que empecé el chat fueron más.md. La máquina que lo compuso lo hizo desde este mismo archivo, y no agregó nada que no estuviera escrito. Si algo de esto es falso, se dice en el mismo sitio, y se corrige con fecha.

\end{document}
